% !PRINTED 748
\weberpage{748}
Jeder Theil der natürlichen Zahlenreihe ist eine abzählbare Menge; denn man braucht ja die Zahlen eines solchen Theiles nur nach ihrer Grösse zu ordnen, um eine zählbare Anordnung zu erhalten.

Daraus aber ergiebt sich, dass jeder Theil einer abzählbaren Menge selbst eine abzählbare Menge ist. Es folgt daraus unter Anderem, dass auch die Menge der \gesperrt{reellen algebraischen} Zahlen abzählbar ist.

\begin{center}
\S.\,204.

\gesperrt{Unzählbare Mengen.}
\end{center}

Wir kommen nun zweitens zu dem Beweise des Satzes, dass es unter den Zahlenmengen auch nicht abzählbare giebt, und wir werden speciell nachweisen:

\begin{center}
\gesperrt{Dass die Gesammtheit aller reellen Zahlen, selbst wenn wir uns auf ein endliches Intervall beschränken, nicht abzählbar ist.}
\end{center}

Wir betrachten zu diesem Zwecke irgend eine abzählbare Menge reeller von einander verschiedener Zahlen, die wir, in eine zählbare Anordnung $\Omega$ gesetzt, so bezeichnen wollen:
\weq{$\Omega$}{\omega_1, \omega_2, \omega_3, \omega_4, \ldots}
(wobei daran zu erinnern ist, dass dies nicht etwa eine Aufeinanderfolge nach der Grösse bedeuten soll).

Der Kürze wegen wollen wir von zwei Elementen der Reihe $\Omega$ das mit kleinerem Index das \gesperrt{frühere}, das mit grösserem das \gesperrt{spätere} nennen.

Wir nehmen nun irgend zwei reelle Zahlen $\alpha, \beta$ an, so dass $\alpha < \beta$ ist, und zeigen, dass es in dem Intervalle $\delta = (\alpha, \beta)$ mindestens eine Zahl giebt, die nicht in der Reihe $\Omega$ vorkommt. Haben wir eine solche Zahl für jedes Intervall nachgewiesen, so giebt es auch deren unendlich viele, da man ja dieselbe Schlussweise auf jeden Theil des Intervalles $(\alpha, \beta)$ anwenden kann.

Zunächst ist klar, dass unsere Behauptung richtig ist, wenn in irgend einem endlichen Theile des Intervalles $\delta$ nur eine endliche Anzahl von Zahlen der Reihe $\Omega$ enthalten ist, und wir können also ohne Weiteres zu dem Falle übergehen, dass in jedem noch so kleinen Theile des Intervalles $\delta$ eine unendliche Menge von Zahlen der Reihe $\Omega$ liegt.
